\chapter{The s-Plane: Poles, Zeros, and Natural Modes}
\label{ch:splane}
\section{One Picture Behind Every View}
So far the same circuit has appeared as an ODE, a state-space model with
eigenvalues, a phasor calculation, and a Bode plot. The \emph{s-plane} is the
single diagram that ties those together. Every pole is a natural mode of the
circuit, every eigenvalue of the state matrix is one of those poles, and the
position of a pole in the complex plane tells you at a glance how fast the mode
decays, how fast it rings, and how sharp the resonance looks on a Bode plot.

This chapter is the hinge between the transient and resonance chapters that
precede it and the filter, matching, and feedback chapters that follow. Nothing
here is a calculation the exam asks you to perform; it is the mental model that
makes the exam's vocabulary---selectivity, bandwidth, ringing, stability---one
idea instead of many.

\begin{controlsbox}
\textbf{Refresher: poles and zeros.} Write a transfer function as a ratio of
polynomials in the complex variable \(s=\sigma+\jj\omega\),
\[
G(s)=\frac{N(s)}{D(s)}.
\]
The \emph{poles} are the roots of \(D(s)\); the \emph{zeros} are the roots of
\(N(s)\). Poles are where \(G\) blows up; zeros are where \(G\) vanishes. For a
state-space model \(\dot{\state}=A\state+Bu\), \(y=C\state+Du\), the poles are
exactly the eigenvalues of \(A\)---the same eigenvalues you compute for stability
and natural modes. The frequency response is the transfer function evaluated on
the imaginary axis, \(G(\jj\omega)\): you slide \(s\) up the \(\jj\omega\) axis
and read off gain and phase.
\end{controlsbox}

\section{A Pole Is a Natural Mode}
The unforced circuit relaxes as a sum of terms \(\ee^{s_i t}\), one per pole
\(s_i\). Writing a pole as \(s_i=\sigma_i+\jj\omega_i\),
\[
\ee^{s_i t}=\ee^{\sigma_i t}\bigl(\cos\omega_i t+\jj\sin\omega_i t\bigr).
\]
The real part \(\sigma_i\) is a growth or decay rate; the imaginary part
\(\omega_i\) is a ringing frequency. This gives the entire stability picture:
\begin{description}
\item[Left half-plane (\(\sigma_i<0\))] the mode decays; the circuit is stable.
\item[Imaginary axis (\(\sigma_i=0\))] the mode neither grows nor decays; a pure
oscillation. This is the knife-edge an oscillator is built to sit on.
\item[Right half-plane (\(\sigma_i>0\))] the mode grows; the circuit is unstable.
\end{description}

For the first-order RC and RL circuits there is a single real pole. The RC
low-pass has
\[
G(s)=\frac{1}{1+sRC},
\qquad
s_{\text{pole}}=-\frac{1}{RC}=-\frac{1}{\tau}.
\]
The pole sits on the negative real axis at minus the reciprocal of the time
constant. A pole far to the left (large \(|\sigma|\), small \(\tau\)) means fast
settling and a high corner frequency; a pole near the origin means a slow,
low-frequency circuit. There is no imaginary part, so a first-order circuit
never rings.

\section{The Second-Order Pole Pair}
The series RLC circuit is where the s-plane earns its keep. From the resonance
chapter, the characteristic polynomial is
\[
LC\,s^2+RC\,s+1=0
\qquad\Longleftrightarrow\qquad
s^2+\frac{R}{L}\,s+\frac{1}{LC}=0.
\]
Introduce the undamped natural frequency and the damping parameters
\[
\wzero=\frac{1}{\sqrt{LC}},
\qquad
\alpha=\frac{R}{2L},
\qquad
\zeta=\frac{\alpha}{\wzero}=\frac{R}{2}\sqrt{\frac{C}{L}},
\]
so that the polynomial takes the standard second-order form
\[
s^2+2\zeta\wzero\,s+\wzero^2=0.
\]
Here \(\zeta\) is the \emph{damping ratio} (dimensionless) and \(\alpha\) is the
damping \emph{rate} (in \si{\per\second}). The poles are
\[
s=-\alpha\pm\sqrt{\alpha^2-\wzero^2}
=-\zeta\wzero\pm\jj\,\wzero\sqrt{1-\zeta^2}
=-\alpha\pm\jj\,\omega_d,
\]
recovering the damped natural frequency \(\omega_d=\sqrt{\wzero^2-\alpha^2}\)
from the resonance chapter as the imaginary part of the pole.

\begin{controlsbox}
\textbf{Refresher: \(\zeta\) and \(Q\) are the same statement.} Control texts
grade a second-order system by its damping ratio \(\zeta\); radio grades a
resonator by its quality factor \(Q\). They are reciprocals up to a factor of
two,
\[
Q=\frac{\wzero}{2\alpha}=\frac{1}{2\zeta}
=\frac{\wzero L}{R},
\]
which is exactly the series-\(Q\) formula from the resonance chapter. Low
damping \(\Leftrightarrow\) high \(Q\) \(\Leftrightarrow\) poles near the
\(\jj\omega\) axis. The three regimes you have been calling overdamped,
critically damped, and underdamped are just \(\zeta>1\), \(\zeta=1\), and
\(\zeta<1\).
\end{controlsbox}

\section{Reading the Geometry}
Two facts make the pole pair readable by eye. First, while the circuit is
underdamped the poles lie on a circle of radius \(\wzero\):
\[
|s|^2=(\zeta\wzero)^2+\bigl(\wzero\sqrt{1-\zeta^2}\bigr)^2=\wzero^2.
\]
The distance of the pole from the origin is the natural frequency. Second, the
angle \(\theta\) measured from the negative real axis satisfies
\[
\cos\theta=\zeta.
\]
So the angle encodes the damping, and therefore \(Q\). A pole hugging the
\(\jj\omega\) axis (\(\theta\to 90^\circ\), \(\zeta\to0\)) is a high-\(Q\),
lightly damped resonance; a pole near the negative real axis is heavily damped.

The real part \(\alpha=\zeta\wzero\) does double duty. It is the envelope decay
rate of the ringing---the transient dies as \(\ee^{-\alpha t}\)---and it also
sets the half-power bandwidth of the resonance,
\[
\text{BW}=2\alpha=\frac{\wzero}{Q}\quad[\si{\radian\per\second}].
\]
\emph{Distance from the imaginary axis is bandwidth.} A pole close to the axis
rings for a long time and shows a narrow, tall peak; a pole far from the axis
settles quickly and shows a broad, low peak. Selectivity in the frequency domain
and ringing in the time domain are the same pole seen two ways.

\section{Root Locus: Watching the Poles Move with \texorpdfstring{$R$}{R}}
Fix \(L\) and \(C\) and sweep the resistance \(R\) from zero upward. Because the
poles solve \(s^2+2\alpha s+\wzero^2=0\) with \(\alpha=R/2L\), they trace a path
in the s-plane as \(R\) grows. This path is a \emph{root locus}.

\begin{figure}[htbp]
\centering
\begin{tikzpicture}[>=Latex,scale=1.15]
% axes
\draw[->,black!70] (-3.6,0) -- (0.9,0) node[right]{\(\sigma\) (real)};
\draw[->,black!70] (0,-2.6) -- (0,2.6) node[above]{\(\jj\omega\) (imag)};
% circle of radius omega0 = 2
\draw[GuideLine,dashed] (0,0) circle (2);
\node[GuideLine] at (1.55,1.75) {\(|s|=\wzero\)};
% R=0 poles on imaginary axis
\filldraw[GuideBlue] (0,2) circle (2.4pt);
\filldraw[GuideBlue] (0,-2) circle (2.4pt);
\node[GuideBlue,right] at (0.08,2.15) {\small \(R=0\)};
% underdamped poles on the circle (angle 150 deg from +real axis)
\filldraw[GuideGreen] (-1.732,1.0) circle (2.4pt);
\filldraw[GuideGreen] (-1.732,-1.0) circle (2.4pt);
% critical damping (double pole) at -omega0
\filldraw[GuideAmber] (-2,0) circle (2.4pt);
\node[GuideAmber,below left] at (-2.0,-0.08) {\small critical};
% overdamped real poles (product = omega0^2 = 4: -1.25 * -3.2 = 4.0)
\filldraw[black!70] (-1.25,0) circle (2.4pt);
\filldraw[black!70] (-3.2,0) circle (2.4pt);
% arrows showing motion as R increases
\draw[->,GuideGreen,thick] (0,1.98) to[bend right=18] (-1.7,1.05);
\draw[->,GuideGreen,thick] (0,-1.98) to[bend left=18] (-1.7,-1.05);
\draw[->,black!70,thick] (-1.9,0.05) to[bend right=25] (-1.3,0.02);
\draw[->,black!70,thick] (-2.1,0.05) to[bend left=25] (-3.15,0.02);
% damping-angle marker
\draw[black!55] (0,0) -- (-1.732,1.0);
\node at (-0.55,0.55) {\small \(\theta\)};
\node[GuideGreen] at (-2.55,1.35) {\small increasing \(R\)};
\end{tikzpicture}
\caption{Root locus of the series RLC poles as \(R\) increases with \(L,C\)
fixed. At \(R=0\) the poles sit on the \(\jj\omega\) axis at \(\pm\jj\wzero\)
(undamped). As \(R\) grows they slide along the circle of radius \(\wzero\)
into the left half-plane, meet at \(-\wzero\) at critical damping
(\(R=2\sqrt{L/C}\)), then split and move apart along the negative real axis
(overdamped). The angle \(\theta\) from the negative real axis satisfies
\(\cos\theta=\zeta\).}
\label{fig:rlc-rootlocus}
\end{figure}

The story in \Cref{fig:rlc-rootlocus} is worth narrating because it unifies
three chapters:
\begin{itemize}
\item \textbf{\(R=0\):} poles on the imaginary axis at \(\pm\jj\wzero\). No loss,
infinite \(Q\), a pure oscillation---the idealization an LC tank approaches.
\item \textbf{Small \(R\) (underdamped):} poles just inside the left half-plane,
high on the circle. High \(Q\), narrow bandwidth, long ringing, tall Bode peak.
\item \textbf{\(R=2\sqrt{L/C}\) (critical):} the pair meets at \(-\wzero\) on the
real axis, \(\zeta=1\). Fastest settling without overshoot.
\item \textbf{Large \(R\) (overdamped):} two real poles; the circuit behaves like
two cascaded first-order lags and no longer rings.
\end{itemize}

\begin{advancedbox}
The product of the two real poles is fixed at \(\wzero^2\) (the constant term of
the monic polynomial), so as \(R\to\infty\) one overdamped pole races toward
\(-\infty\) while the other creeps back toward the origin at
\(s\approx-\wzero^2/(2\alpha)=-1/(RC)\). The slow pole becoming the dominant RC
lag is the same limit you would expect physically: a huge series resistor makes
the capacitor charge through \(R\) alone.
\end{advancedbox}

\section{How Pole Position Shapes the Bode Plot}
Because the frequency response is \(G(\jj\omega)\)---the transfer function
sampled along the imaginary axis---the poles you just placed control the Bode
plot directly.

\begin{controlsbox}
\textbf{Refresher: poles, zeros, and Bode slopes.} Each pole bends the magnitude
slope down by \(\SI{20}{\decibel}\) per decade near its corner frequency and adds
up to \(-90^\circ\) of phase; each zero bends it up by
\(+\SI{20}{\decibel}\) per decade and adds up to \(+90^\circ\). The corner sits
at the frequency equal to the pole's or zero's distance from the origin. This is
why the first-order rules in the frequency-response chapter worked: a ``pole'' is
just a corner in the asymptotic plot.
\end{controlsbox}

When an underdamped pole pair sits close to the \(\jj\omega\) axis, sliding
\(\omega\) past the pole's height \(\omega_d\approx\wzero\) brings the evaluation
point \(\jj\omega\) very near the pole, where \(|G|\) grows large. That near-miss
is the resonant peak. Its height is governed by how close the pole is to the
axis, i.e.\ by \(Q\): peak magnitude scales with \(Q\), and the \(-\SI{3}{dB}\)
width of the peak equals the bandwidth \(2\alpha\). High \(Q\), pole near the
axis, tall narrow peak; low \(Q\), pole pulled left, gentle bump. A
\emph{zero} on or near the axis does the opposite---it pulls \(|G|\) down toward
a notch, which is exactly how a trap or notch filter rejects one frequency.

\section{Zeros in Radio Circuits}
Zeros are less discussed than poles but they are how circuits \emph{reject}
energy. A series LC branch to ground has zero impedance at its resonance, shorting
that frequency to ground and placing a transmission zero on the \(\jj\omega\)
axis at \(\wzero=1/\sqrt{LC}\); this is a trap. A parallel LC branch in series
with the signal path does the dual: it becomes an open circuit at resonance,
again blocking one frequency. Wherever the exam describes a trap, a notch, or a
wave trap rejecting a specific frequency, a zero has been parked at that
frequency.

\section{Worked Example: Locating the Poles of a Tuned Circuit}
Reuse the resonance-chapter values \(L=\SI{10}{\micro\henry}\),
\(C=\SI{220}{\pico\farad}\), \(R=\SI{5}{\ohm}\). Then
\[
\wzero=\frac{1}{\sqrt{LC}}\approx\SI{2.13e7}{\radian\per\second}
\quad(f_0\approx\SI{3.39}{\mega\hertz}),
\qquad
\alpha=\frac{R}{2L}=\frac{5}{2(10\times10^{-6})}=\SI{2.5e5}{\radian\per\second}.
\]
The damping ratio and quality factor are
\[
\zeta=\frac{\alpha}{\wzero}\approx 0.0117,
\qquad
Q=\frac{1}{2\zeta}\approx 42.7,
\]
matching the \(Q\approx42.6\) found earlier. The poles sit at
\[
s=-\alpha\pm\jj\,\omega_d\approx-2.5\times10^{5}\pm\jj\,2.13\times10^{7}
\ \si{\radian\per\second}.
\]
The real part is about \(1/85\) of the imaginary part, so the poles are pressed
right against the \(\jj\omega\) axis: a very high-\(Q\), lightly damped,
sharply peaked circuit whose ringing decays only after roughly \(Q\) cycles.
Everything the resonance chapter said in the language of \(Q\) and bandwidth is
visible here as a pole location.

\begin{exambox}
The exam does not ask you to plot poles or compute \(\zeta\). It asks about
resonant frequency, \(Q\), bandwidth, selectivity, ringing, oscillator
stability, and filter sharpness---all of which are the answer to a single
question: \emph{where are the poles?} Near the \(\jj\omega\) axis means sharp,
selective, and long-ringing; deep in the left half-plane means broad and
well-damped; on the axis means a sustained oscillation.
\end{exambox}
